% p.72 — العوامل ومتعلقاتها (تتمة)
% Layout: twocol. rightcol: the last example line of entry (106) بدل from p.71, 107–109,
% sub-heading حروف العطف عشرة, 110–119; leftcol: the rule paragraph on غير المنصرف, 120–128,
% sub-heading اسماء الموصولات, 129–131. Rule paragraph is \par-separated text between the tables.
% Note: reledpar sets the two columns line by line (line i of the right column beside line i of
% the left one, each pair as tall as its taller member), so a column made of tall sarf "lines"
% padded the other column and the page overflowed. Each column's whole content is therefore
% wrapped in one rightbox/leftbox (minipages inside the reledpar columns): the columns then
% flow independently and the page fits; the apparatus numbers each column as a single line.
\mspage{72}
\begin{twocol}
\begin{rightbox}
\begin{sarf}{|c|c|}
 & اَكَلْتُ الرَّغِيْفَ نِصْفَهُ - رَغِبْتُ فِيْ زَيْدٍ عِلْمِهِ\\
(107) عَطْفُ الْبَيَانِ & اَظْهَرَ الْاِسْلَامَ اَبُوْ حَفْصٍ عُمَرُ\\
 & اَحْبَبْتُ بَيْتَ اللهِ الْكَعْبَةَ\\
 & مَرَرْتُ بِعَبْدِ اللهِ ابْنِ عُمَرَ\\
(108) تَأْكِيْدٌ & سَجَدَ الْمَلَائِكَةُ كُلُّهُمْ اَجْمَعُوْنَ\\
 & رَاَيْتُ زَيْدًا عَيْنَهُ\\
 & مَرَرْتُ بِزَيْدٍ نَفْسِهِ\\
(109) عَطْفٌ بِالْحَرْفِ & جَاءَنِيْ زَيْدٌ وَعَمْرٌو\\
 & رَاَيْتُ زَيْدًا اَوْ عَمْرًا\\
 & مَرَرْتُ بِزَيْدٍ ثُمَّ بِعَمْرٍو\\
\end{sarf}
{\subbab{حُرُوْفُ الْعَطْفِ عَشَرَةٌ}}
\begin{sarf}{|c|c|}
(110) اَلْوَاوُ & جَاءَنِيْ زَيْدٌ وَعَمْرٌو\\
(111) اَوْ & رَاَيْتُ زَيْدًا اَوْ بَكْرًا\\
(112) ثُمَّ & جَاءَنِيْ بَكْرٌ ثُمَّ خَالِدٌ\\
(113) فَاءٌ & \unclear{تَخَرَّجَ} زَيْدٌ فَعَمْرٌو\\ % TODO: verb unclear (تَخَرَّجَ / تَحَرَّجَ)
(114) حَتّٰى & مَاتَ النَّاسُ حَتَّى الْاَنْبِيَاءُ\\
(115) اِمَّا & اَلْعَدَدُ اِمَّا زَوْجٌ وَاِمَّا فَرْدٌ\\
(116) اَمْ & اَعِنْدَكَ زَيْدٌ اَمْ عَمْرٌو\\
(117) لَا & رَاَيْتُ زَيْدًا لَا خَالِدًا\\
(118) بَلْ & مَا جَاءَ بَكْرٌ بَلْ خَالِدٌ\\
(119) لٰكِنْ & مَا جَاءَنِيْ زَيْدٌ لٰكِنْ خَالِدٌ جَاءَ\\
\end{sarf}
\end{rightbox}
\begin{leftbox}
اِعْلَمْ اَنَّ الْاِسْمَ اِذَا وُجِدَ فِيْهِ عِلَّتَانِ مِنَ التِّسْعِ يَكُوْنُ الْاِسْمُ غَيْرَ مُنْصَرِفٍ وَمَعْنٰى عَدَمِ الصَّرْفِ اَنْ لَا يَظْهَرَ فِيْهِ الْكَسْرَةُ وَلَا يَحْمِلُ التَّنْوِيْنَ وَهِيَ تِسْعَةٌ
\par
\begin{sarf}{|c|c|}
(120) عَدْلٌ & جَاءَنِيْ عُمَرُ\\
(121) وَصْفٌ & رَاَيْتُ رَجُلًا اَسْوَدَ\\
(122) تَأْنِيْثٌ & هٰذَا طَلْحَةُ رَاَيْتُ زَيْنَبَ\\
 & مَرَرْتُ بِحَمْرَاءَ\\
(123) عُجْمَةٌ & مَرَرْتُ بِاِبْرَاهِيْمَ\\
(124) جَمْعٌ & دَخَلْتُ الْمَسَاجِدَ وَاَطْفَأْتُ\\
 & الْمَصَابِيْحَ وَضَرَبْتُ الدَّوَابَّ\\
(125) تَرْكِيْبٌ & جَاءَنِيْ بَعْلَبَكُّ\\
(126) اَلْاَلِفُ وَالنُّوْنُ & صُمْتُ رَمَضَانَ هٰذَا\\ % the two example lines are joined by a brace { to the two-line headword
الزَّائِدَتَانِ & شَعْبَانُ مَرَرْتُ بِعُثْمَانَ\\
(127) وَزْنُ فِعْلٍ & هٰذَا اَحْمَدُ وَرَاَيْتُ \unclear{تَرْحَبَ}\\ % TODO: last word smudged (perhaps تَغْلِبَ)
(128) عَلَمٌ & ضَرَبْتُ زُبَيْرًا وَمَرَرْتُ بِعُمَرَ\\
\end{sarf}
{\subbab{اَسْمَاءُ الْمَوْصُوْلَاتِ}}
\begin{sarf}{|c|c|}
(129) اَلَّذِيْ & جَاءَ رَجُلٌ اَلَّذِيْ ضَرَبْتُهُ\\
(130) اَلَّذَانِ & جَاءَ رَجُلَانِ اَلَّذَانِ ضَرَبْتُهُمَا\\
 & رَاَيْتُ رَجُلَيْنِ اَلَّذَيْنِ ضَرَبْتُهُمَا\\
(131) اَلَّذِيْنَ & ضَرَبْتُ رِجَالًا اَلَّذِيْنَ ضَرَبُوْكَ\\
\end{sarf}
\end{leftbox}
\end{twocol}
